\section{Numerical operator and implementation}
\label{app:appendix_math}

\subsection{Finite Newton--Schulz map}

The executed map uses ten polynomial iterations in mixed precision, so it is not an exact SVD polar factor. For $Z=LUR$, we first compute $X_0=Z/(\|Z\|_F+10^{-5})$ in FP32. Tall matrices are transposed so each Gram has dimension $\min(m,n)$. For $k=0,\ldots,9$,
\begin{equation}
 A_k=X_k\star X_k^{\Tr},\qquad
 X_{k+1}=3.4445X_k+[-4.7750A_k+2.0315(A_k\star A_k)]\star X_k,
 \label{eq:ns10}
\end{equation}
where $P\star Q$ rounds both operands to BF16 and accumulates the matrix product in FP32. Scalar combinations stay in FP32, and tall results are transposed back. The exact identity in \Cref{eq:som_rotation} is a reference for understanding this operator, not an exact description of every finite-iteration update.

The recorded sensitivity check perturbs each FP32 input element by one ULP. The perturbation grows through iterations 3--10 to a $5.2\times10^{-4}$ parameter-delta shift under this rounding path; perturbing outer factors gives a $1.9\times10^{-9}$ shift. These checks motivate the unified execution path and Null-WT comparison. Deterministic runs use \texttt{CUBLAS\_WORKSPACE\_CONFIG=:4096:8}.

\subsection{Dense products and cotangent capture}

Dense products outside the Newton--Schulz map use an FP16-split Tensor-Core implementation. With TF32 disabled, this avoids the slow native FP32 path on A100. For a tensor $X$, the split uses $s=2^{13-\lfloor\log_2\max|X|\rfloor}$, $X_{\rm hi}=\operatorname{half}(sX)$, and $X_{\rm lo}=\operatorname{half}(4096(sX-\operatorname{float}(X_{\rm hi})))$. The recorded product audit finds relative error below $3\times10^{-7}$ against FP64 references. This is an empirical arithmetic check for the tested matrices.

Backward capture receives loss-scaled cotangents. For $G_{\rm acc}$ microbatches, $B_{\rm micro}$ sequences, and 4,095 scored targets per sequence, autograd returns the summed-loss derivative divided by $B_{\rm micro}\,4095\,G_{\rm acc}$. Capture applies a fixed $2^{16}$ scale before BF16 storage and FP32 accumulation to reduce underflow. Raw sums and sample counts are pooled across ranks before the EMA update. Trace normalization cancels a common constant scale in exact arithmetic; finite-precision capture can still introduce rounding error.

The core methods share MuonT's weight-decay rule $Q=W\hat R^2$. Output-only ablations use $Q=\hat L^2W$, and Muon uses $Q=W$. Thus, output-only replacement ablations also change the decay operator and cannot alone isolate the role of input whitening. The separate related-method comparison uses isotropic decay for every arm.

\section{Moment audit and geometric diagnostics}
\label{app:appendix_mech}

\subsection{Definitions and replay checks}
\label{app:diagnostic_definitions}

The audit asks three distinct questions: how much moment structure a diagonal approximation retains, how closely that approximation reproduces the dense inverse root, and how much the leading eigenspace changes during training. These questions require separate diagnostics. Throughout this subsection, $F_t\in\R^{d\times d}$ denotes one symmetric positive-semidefinite EMA second-moment matrix from \cref{eq:moments}: $C_{g(p),t}$ for an input activation or $B_{p,t}$ for an output cotangent, at a fixed layer and projection. We assume $\tr(F_t)>0$ and use $\tilde F_t=dF_t/\tr(F_t)$, whose mean diagonal entry and mean eigenvalue are one. All moments are uncentered. We omit $t$ for diagnostics evaluated at a single checkpoint.

We write $\diag(M)$ for the vector of diagonal entries, $\Diag(a)$ for the diagonal matrix with entries $a$, and $\|M\|_F^2=\sum_{i,j}M_{ij}^2$ for squared Frobenius norm. Static columns in \Cref{tab:diagnosis} are medians of per-layer scores across 32 layers at step 2,500; temporal drift is the mean of 32 per-layer comparisons between steps 2,250 and 2,500. Each optimizer trajectory is evaluated separately.

Diagonal anisotropy share measures where the departure from isotropy lies in the model's coordinate basis. Let $D=\Diag(\diag(\tilde F))$. Since diagonal and off-diagonal matrices are orthogonal under the Frobenius inner product,
\[
 \|\tilde F-I_d\|_F^2
 =\underbrace{\|D-I_d\|_F^2}_{E_{\rm diag}}
 +\underbrace{\|\tilde F-D\|_F^2}_{E_{\rm off}}.
\]
The best diagonal approximation to $\tilde F$ in this norm is $D$. Replacing identity with $D$ therefore removes exactly $E_{\rm diag}$ of the squared approximation error. This gives $\mathrm{DiagAniso}=E_{\rm diag}/(E_{\rm diag}+E_{\rm off})$ in \cref{eq:diag_aniso}, and explains its use in Figures~\ref{fig:overview}(a) and~\ref{fig:structural_audit}(a). A value of 0.8 means that 80\% of the departure from identity is diagonal. It does not specify the magnitude of that departure. The ratio is mathematically undefined for an exactly isotropic moment, where both contributions vanish; the diagnostic code returns zero in that case as a numerical convention.

Diagonal coefficient of variation (CV) supplies the missing measure of coordinate unevenness:
\[
 \mathrm{CV}_{\rm diag}(F)
 =\frac{\sqrt{d^{-1}\sum_i(F_{ii}-\bar f)^2}}{\bar f}
 =\sqrt{\frac{E_{\rm diag}}{d}},
 \qquad \bar f=\frac{\tr(F)}{d}.
\]
It is zero when all coordinate second moments are equal, and increases with their relative spread. A non-isotropic diagonal matrix has $\mathrm{DiagAniso}=1$ even when its CV is arbitrarily small. Reading the two columns together separates the share of diagonal structure from its strength.

Off-diagonal correlation RMS measures typical pairwise coupling after removing each coordinate's scale. For positive diagonal entries and $d>1$, define
\[
 \rho_{ij}=\frac{F_{ij}}{\sqrt{F_{ii}F_{jj}}},\qquad
 \rho_{\rm rms}=\sqrt{\frac{1}{d(d-1)}\sum_{i\ne j}\rho_{ij}^2}.
\]
For a Gram matrix, $\rho_{ij}$ is the cosine similarity between the sampled values of coordinates $i$ and $j$, with the same weights used to form the moment. Squaring prevents positive and negative cross-moments from canceling. Because the samples are not mean-subtracted, this is an uncentered correlation: nonzero coordinate means can contribute even when centered covariance is zero. A zero diagonal in a positive-semidefinite moment implies a zero row and column; such pairs contribute zero, with the same $d(d-1)$ denominator. It differs from $1-\mathrm{DiagAniso}$: it averages individually normalized pairs, whereas the latter measures their aggregate contribution relative to total anisotropy. In a wide matrix, many small correlations can still contribute appreciable off-diagonal energy.

Inverse-root fidelity asks how closely the diagonal factor reproduces the dense inverse root. We use the normalization and damping of the inner factors in \cref{eq:dense_roots,eq:diag_roots}:
\[
 H=(\tilde F+\epsilon I_d)^{-1/2},\qquad
 H_{\rm diag}=\Diag\!\left((\diag(\tilde F)+\epsilon\mathbf1)^{-1/2}\right),
 \qquad\epsilon=10^{-3},
\]
where the power inside $\Diag$ is elementwise. We report the following identity-relative score as root $R^2$:
\begin{equation}
 R^2=1-\frac{\|H-H_{\rm diag}\|_F^2}{\|H-I_d\|_F^2}.
 \label{eq:r2_root}
\end{equation}
The denominator is the squared error from using identity, our no-preconditioning reference; the numerator is the error after using the diagonal factor. Thus $R^2=1$ means exact agreement with the dense root, $R^2=0$ means the same error as identity, and $R^2<0$ means greater error. For example, $R^2=0.8$ means an 80\% reduction in squared error relative to identity. This is a diagnostic defined for this comparison, with the normalization of a relative squared-error reduction; it is not a regression fit or a fraction of recovered eigenvectors. Its interpretation becomes sensitive when the reference error is very small. For example, $\tilde F=I_d$ gives $H=H_{\rm diag}=(1+\epsilon)^{-1/2}I_d$ and $R^2=1$, although the only change is scalar damping. Perfect fidelity alone does not imply a substantial preconditioning effect.

The Frobenius norm has an operator interpretation: for a probe vector $z$ with $\E[zz^{\Tr}]=I_d$, $\E\|(H-H_{\rm diag})z\|_2^2=\|H-H_{\rm diag}\|_F^2$. It measures average squared disagreement on isotropic probes. Actual optimizer updates need not have that distribution, so this score does not determine training-loss improvement. A separate root comparison is necessary because inverse square roots amplify small-eigenvalue directions: a high diagonal share of the moment alone does not guarantee high fidelity of its inverse root.

Subspace drift measures the temporal stability of the directions with the largest second moments. For a unit direction $a$, $a^{\Tr}F_ta$ is the EMA of the squared activation or cotangent projected onto $a$. The leading eigenvectors select orthogonal directions with the largest such values. For each checkpoint, order the eigenvalues of the same $F_t$ as $\lambda_{1,t}\ge\cdots\ge\lambda_{d,t}$ and choose orthonormal eigenvectors satisfying $F_tv_{i,t}=\lambda_{i,t}v_{i,t}$. Here $k=16$ and all audited dimensions satisfy $d\ge512$. Form
\[
 V_{k,t}=[v_{1,t},\ldots,v_{k,t}]\in\R^{d\times k},\qquad
 V_{k,t}^{\Tr}V_{k,t}=I_k.
\]
Each column is a direction across the $d$ activation or cotangent coordinates. For example, a width-512 moment gives a $512\times16$ matrix. The eigenvectors are computed from the full diagnostic moment, including for expansion outputs whose active optimizer uses only its diagonal. Trace normalization and adding $\epsilon I_d$ preserve these eigenspaces. The temporal score is
\begin{equation}
 \mathrm{Drift}_{16}(t)=1-\frac1{16}\|V_{16,t}^{\Tr}V_{16,t-250}\|_F^2.
 \label{eq:subspace_drift}
\end{equation}

To see why this measures subspace change, let $V=V_{k,t}$, $U=V_{k,t-250}$, and $P=UU^{\Tr}$, the orthogonal projector onto the earlier subspace. For a current unit vector $v_j$, $Pv_j$ is its part inside that subspace and $(I_d-P)v_j$ is the residual outside it. Averaging the squared residual gives
\[
 \frac1k\sum_{j=1}^k\|(I_d-P)v_j\|_2^2
 =\frac1k\|(I_d-P)V\|_F^2
 =1-\frac1k\|U^{\Tr}V\|_F^2.
\]
Thus \cref{eq:subspace_drift} is the mean fraction of each current unit basis vector's squared norm lying outside the earlier span. It equals zero for identical subspaces and one for orthogonal subspaces. If 15 directions are shared and the remaining direction becomes orthogonal to the old span, the score is $1/16$. Sign flips, permutations, and orthogonal changes of basis within either span leave the score unchanged.

This construction is a normalized squared subspace distance expressed through principal angles~\citep{bjorck1973angles,yu2015davis}. If $\theta_1,\ldots,\theta_k\in[0,\pi/2]$ are the principal angles, defined by $\cos\theta_j=\sigma_j(U^{\Tr}V)$, then
\[
 \mathrm{Drift}_k
 =\frac1k\sum_{j=1}^k\sin^2\theta_j
 =\frac{\|VV^{\Tr}-UU^{\Tr}\|_F^2}{2k}.
\]
For $k=1$, it reduces to $1-(u^{\Tr}v)^2=\sin^2\theta$, the squared residual from projecting one unit direction onto another. The choices $k=16$ and 250 updates specify the diagnostic resolution and interval; they are not required by the distance definition. Figure~\ref{fig:structural_audit}(b) uses $t=2500$ and compares the same moment in the same layer along each trajectory.

Drift helps assess whether a cached leading eigenspace remains representative as training proceeds. It gives equal weight to the selected directions and does not measure changes in their eigenvalues, inverse-root error, or a per-step rotation rate. Directions can enter or leave the leading set through eigenvalue crossings. A small gap between $\lambda_{16}$ and $\lambda_{17}$ can also make that set sensitive to small perturbations; subspace perturbation bounds explicitly depend on eigenvalue separation~\citep{yu2015davis}. An exact tie across this boundary makes the selected 16-dimensional subspace nonunique. The checkpoint-audit scores use the eigensolver's leading vectors without an eigenvalue-gap filter; \Cref{sec:output_drift} adds a gap-filtered sensitivity check for the trajectory comparisons. The audit therefore uses drift as a descriptive temporal diagnostic and reports observed root-solver fallbacks separately, rather than inferring a required repair rate from drift alone.

Effective rank summarizes how broadly the moment's spectral mass is distributed. Using the nonnegative eigenvalues of $F$, we follow the entropy-based definition of \citet{roy2007effective}:
\[
 p_i=\frac{\lambda_i}{\sum_j\lambda_j},\qquad
 \operatorname{erank}(F)=\exp\!\left(-\sum_i p_i\log p_i\right),
 \qquad 0\log0=0.
\]
The table reports $100\operatorname{erank}(F)/d$. If exactly $r$ eigenvalues are equal and positive and the rest are zero, the effective rank is $r$; concentration in one direction gives one, and an isotropic moment gives $d$. This measures spectral concentration even when the ordinary rank is full. It is invariant to orthogonal changes of coordinates, so it complements the coordinate-dependent CV and diagonal-share columns.

The replay checks distinguish accumulator consistency from reproduction of an uninterrupted training trajectory. Muon and WT-MLP-Lag restart from step 2,000, with diagnostics saved at steps 2,250 and 2,500. Passive moments start at zero and are divided by the accumulated EMA weight $1-\beta^n$ after $n$ updates. For identical subsequent samples, a discrepancy in the uncorrected initial state decays as $\beta^n\Delta F_0$; with $\beta=0.938$, the coefficient is $1.12\times10^{-7}$ after 250 updates. This describes forgetting of initialization, without bounding sampling error. Muon collects all 11 moments passively. WT-MLP-Lag uses its active dense moments and collects full expansion-output moments passively, without applying the latter to the update.

Muon reproduces the recorded losses and checkpoint parameters bitwise. WT-MLP-Lag instead shows a maximum absolute training-loss difference of 0.00543 nats over the replay window and a maximum parameter-entry difference of 0.333 at step 2,500 relative to its original checkpoint. Its reported geometry therefore describes the resumed trajectory. Within that replay, let $b_r^{\rm obs}$ concatenate the diagonals of the observed full gate/up moments owned by rank $r$, and let $b_r^{\rm act}$ concatenate the corresponding stored active diagonal accumulators, without additional rescaling. The passive correction factor at this checkpoint differs from one by $1.26\times10^{-14}$. At step 2,500,
\[
 \max_{r=0,\ldots,7}
 \frac{\|b_r^{\rm obs}-b_r^{\rm act}\|_2}{\|b_r^{\rm act}\|_2}
 =4.04\times10^{-5}.
\]
For the stored original-checkpoint diagonals $b_r^{\rm orig}$, the comparison $\max_r\|b_r^{\rm obs}-b_r^{\rm orig}\|_2/\|b_r^{\rm orig}\|_2$ instead gives 0.558. The small within-replay error establishes aggregate diagonal agreement at step 2,500. It does not validate every individual moment, the off-diagonal entries, or reproduction of the original WT trajectory. Neither replay check establishes that the observed moments equal population covariances or true curvature.

\begin{table}[tbp]\centering\footnotesize
\setlength{\tabcolsep}{8.0pt}
\caption{\textbf{Complete 352-matrix moment audit across three 130M trajectories.} Effective rank is a percentage of width $d$; CV is diagonal coefficient of variation; $\rho_{\rm rms}$ is off-diagonal correlation RMS; DA is diagonal anisotropy share (\%); root $R^2$ is diagonal inverse-root explained variance (\%). Static metrics are 32-layer medians at step 2,500; drift is the 32-layer mean over steps 2,250 to 2,500.}\label{tab:diagnosis}
\begin{tabularx}{\textwidth}{@{}>{\raggedright\arraybackslash}Xrrrrrrr@{}}
\toprule
Moment & $d$ & Rank (\%) & CV & $\rho_{\rm rms}$ & DA (\%) & Root $R^2$ & Drift \\
\midrule
\addlinespace[2pt]\multicolumn{8}{@{}l}{\textit{Muon}}\\
Input: attention & 512 & 30.0 & 0.63 & 0.1465 & 1.4 & 8.2 & 0.041 \\
Input: o & 512 & 37.9 & 0.77 & 0.0968 & 9.8 & 30.1 & 0.107 \\
Input: MLP & 512 & 33.8 & 0.48 & 0.1410 & 1.0 & 4.9 & 0.046 \\
Input: down & 2,048 & 67.1 & 0.71 & 0.0283 & 20.3 & 91.3 & 0.186 \\
Output: q & 512 & 38.5 & 1.04 & 0.0615 & 16.4 & 64.8 & 0.150 \\
Output: k & 512 & 24.4 & 1.45 & 0.0746 & 15.5 & 77.6 & 0.178 \\
Output: v & 512 & 47.0 & 0.63 & 0.0600 & 11.5 & 38.6 & 0.664 \\
Output: o & 512 & 83.3 & 0.09 & 0.0307 & 1.0 & 3.6 & 0.164 \\
Output: down & 512 & 84.3 & 0.08 & 0.0315 & 0.9 & 2.3 & 0.123 \\
Output: gate & 2,048 & 88.2 & 0.55 & 0.0055 & 80.6 & 98.8 & 0.292 \\
Output: up & 2,048 & 83.8 & 0.68 & 0.0068 & 77.9 & 98.2 & 0.251 \\
\addlinespace[4pt]\multicolumn{8}{@{}l}{\textit{MuonT}}\\
Input: attention & 512 & 17.9 & 0.84 & 0.2031 & 1.0 & 10.6 & 0.021 \\
Input: o & 512 & 20.8 & 0.75 & 0.2120 & 1.8 & 22.5 & 0.094 \\
Input: MLP & 512 & 21.6 & 0.64 & 0.1959 & 0.8 & 6.0 & 0.021 \\
Input: down & 2,048 & 29.1 & 2.92 & 0.0346 & 43.4 & 79.4 & 0.106 \\
Output: q & 512 & 38.5 & 0.97 & 0.0641 & 18.0 & 62.4 & 0.140 \\
Output: k & 512 & 23.7 & 1.50 & 0.0774 & 14.4 & 71.8 & 0.165 \\
Output: v & 512 & 50.4 & 0.58 & 0.0560 & 12.2 & 31.2 & 0.662 \\
Output: o & 512 & 85.3 & 0.06 & 0.0301 & 0.8 & 1.2 & 0.134 \\
Output: down & 512 & 84.6 & 0.07 & 0.0312 & 0.8 & 1.2 & 0.107 \\
Output: gate & 2,048 & 87.1 & 0.59 & 0.0051 & 82.0 & 82.2 & 0.240 \\
Output: up & 2,048 & 82.5 & 0.75 & 0.0073 & 78.1 & 75.9 & 0.215 \\
\addlinespace[4pt]\multicolumn{8}{@{}l}{\textit{WT-MLP-Lag}}\\
Input: attention & 512 & 18.0 & 0.76 & 0.2066 & 0.8 & 10.1 & 0.023 \\
Input: o & 512 & 23.4 & 0.77 & 0.2001 & 1.8 & 26.5 & 0.108 \\
Input: MLP & 512 & 19.5 & 0.62 & 0.2024 & 0.7 & 6.7 & 0.025 \\
Input: down & 2,048 & 35.6 & 2.33 & 0.0359 & 34.7 & 78.7 & 0.087 \\
Output: q & 512 & 29.3 & 1.14 & 0.0732 & 13.3 & 66.0 & 0.110 \\
Output: k & 512 & 13.0 & 1.87 & 0.0779 & 14.7 & 78.6 & 0.147 \\
Output: v & 512 & 41.5 & 0.57 & 0.0749 & 6.8 & 19.9 & 0.586 \\
Output: o & 512 & 78.2 & 0.10 & 0.0393 & 0.8 & 2.3 & 0.103 \\
Output: down & 512 & 77.9 & 0.09 & 0.0403 & 0.8 & 1.4 & 0.072 \\
Output: gate & 2,048 & 72.8 & 1.32 & 0.0048 & 89.5 & 94.9 & 0.159 \\
Output: up & 2,048 & 47.3 & 2.83 & 0.0082 & 90.8 & 91.9 & 0.130 \\
\bottomrule\end{tabularx}
\end{table}


In this audit, expansion-output $\rho_{\rm rms}$ ranges from 0.0048 to 0.0082 across the three trajectories, compared with roughly 0.03--0.04 on residual outputs. Together with diagonal anisotropy share and inverse-root fidelity, these values describe the accuracy of the expansion-output approximation under the reported moment and isotropic-probe diagnostics; they do not establish accuracy on every layer or actual update. The correlations are nonzero. Attributing them to sampling noise would require a separate analysis of token dependence and the EMA's effective sample size.

\subsection{Polar rotation and its assumptions}

For full row rank $Z_0=S_0P_0$, the row-polar factor is $(Z_0Z_0^{\Tr})^{-1/2}Z_0$. Substituting $LZ_0$ gives \Cref{eq:som_rotation}. Since $L$ and $S_0$ are positive definite, the square factor $LS_0$ has positive determinant; its polar factor is orthogonal with determinant one. If the two matrices commute, their product is symmetric positive definite and this rotation is identity. The derivation requires $m\le n$ and does not directly apply to the tall gate/up matrices.

At a row-orthogonal $P_0$, an infinitesimal left rotation has the form $KP_0$ with $K^{\Tr}=-K$. An infinitesimal symmetric outer change is $HP_0$ with $H^{\Tr}=H$. Their Frobenius inner product is $\tr(HK^{\Tr})=0$. Finite changes, input outer factors, and finite-iteration polar errors can all change this separation in the executed update (\Cref{tab:mechanism_rotation,tab:rotation_detail}).

\begin{table}[tbp]\centering\footnotesize
\setlength{\tabcolsep}{3.2pt}
\caption{\textbf{Complete 130M inner/outer decomposition and spectral alignment probes (32-layer means, step 2,500).} Inner SVD and Inner $\Nten$ measure left-factor rotation under exact polar decomposition versus ten-step Newton--Schulz; Outer $\Nten$ isolates the post-polar $-1/8$ factor; Cosine compares the inner and outer update deltas (near zero on $m\le n$ projections, confirming orthogonal tangent/normal action). \texttt{bot80} win is the fraction of trailing 80\% $B_t$ modes gaining row alignment; $\mathrm{erank}$ compares pre-polar effective rank ($Z_0\to L Z_0$); $\mathcal Q_{\rm WT}/\mathcal Q_{\rm MT}$ and $\Phi_{\rm OD}/\Phi_{\rm MT}$ complete the quadratic surrogate ratios.}\label{tab:rotation_detail}
\begin{tabularx}{\textwidth}{@{}>{\raggedright\arraybackslash}Xrrrrrrrrr@{}}
\toprule
Output & Inner SVD & Inner $\Nten$ & Outer $\Nten$ & Cosine & Hijack & \texttt{bot80} win & $\mathrm{erank}$ & $\mathcal Q_{\rm WT}/\mathcal Q_{\rm MT}$ & $\Phi_{\rm OD}/\Phi_{\rm MT}$ \\
\midrule
q & $13.4^\circ$ & $21.7^\circ$ & $7.5^\circ$ & $+0.016$ & $67.8\to2.8\%$ & 70.7\% & $65.5\to78.5\%$ & 0.667 & 1.056 \\
k & $14.2^\circ$ & $21.7^\circ$ & $7.7^\circ$ & $+0.027$ & $64.2\to3.7\%$ & 78.6\% & $58.3\to76.3\%$ & 0.607 & 1.064 \\
v & $16.8^\circ$ & $23.1^\circ$ & $8.3^\circ$ & $+0.000$ & $19.6\to1.1\%$ & 69.0\% & $64.9\to68.0\%$ & 0.663 & 1.098 \\
o & $6.9^\circ$ & $15.1^\circ$ & $3.6^\circ$ & $+0.013$ & $30.4\to3.5\%$ & 71.4\% & $74.5\to77.6\%$ & 0.934 & 1.001 \\
down & $2.8^\circ$ & $15.3^\circ$ & $3.5^\circ$ & $-0.002$ & $64.5\to2.8\%$ & 82.4\% & $92.7\to95.9\%$ & 0.895 & 0.998 \\
gate & $10.1^\circ$ & $14.1^\circ$ & $3.1^\circ$ & $+0.579$ & $22.7\to3.1\%$ & 72.7\% & $94.5\to95.3\%$ & 0.767 & 1.292 \\
up & $10.3^\circ$ & $14.7^\circ$ & $3.4^\circ$ & $+0.567$ & $33.7\to4.7\%$ & 70.7\% & $94.1\to95.2\%$ & 0.765 & 1.249 \\
\bottomrule\end{tabularx}
\end{table}


The 130M counterfactual spectral probes use the saved MuonT/null-dense trajectory at step 2,500. The 520M probes use 16 layers from the saved WT-MLP-Lag trajectory at step 5,120, holding that checkpoint's momentum and input factors fixed while changing the output operator. Thus, the input-only counterfactual at 520M is not a separately trained MuonT checkpoint. These measurements isolate operator action within each state and should not be interpreted as a controlled comparison of state distributions across model sizes.

For each output eigenvector $u_i$, the row-alignment score is the cosine between $u_i^{\Tr}\Nten(LZ_0)$ and $u_i^{\Tr}Z_0$, compared with the input-only polar map. The reported bulk-mode fraction uses the trailing 80\% of output eigenvalues. The top-16 overlap is $\|U_{Z,16}^{\Tr}U_{B,16}\|_F^2/16$. Both are descriptive directional statistics.

\subsection{Delocalization reference and quadratic proxy}

For a unit vector $u$, define $\mathrm{IPR}(u)=d\sum_i u_i^4$. It equals $d$ for a coordinate axis and has expectation $3d/(d+2)$ for a Haar-random direction. For a Haar orthogonal eigenbasis $U$ and fixed eigenvalues with mean one and variance $\sigma_\lambda^2$, the diagonal entries of $U\Lambda U^{\Tr}$ have variance $2\sigma_\lambda^2/(d+2)$. The denominator of diagonal anisotropy is $d\sigma_\lambda^2$, yielding \Cref{eq:haar_scaling} when $\sigma_\lambda^2>0$. Similar IPR values in a learned basis establish spread across coordinates, not the full independence and distributional assumptions of the Haar model.

The quadratic score uses trace-normalized moments, $\hat D=D/\|D\|_F$, and $\mathcal Q_{B,C}(\hat D)=\tr(\tilde B\hat D\tilde C\hat D^{\Tr})$. We report $\Phi_{B,C}(\hat D)=\langle U,\hat D\rangle^2/\mathcal Q_{B,C}(\hat D)$ as a local surrogate for alignment relative to moment-weighted perturbation. It uses momentum instead of the loss gradient and observed-label second moments instead of verified curvature. The squared numerator also discards the sign of alignment. A larger value alone therefore cannot certify descent or generalization.

\FloatBarrier
\section{Learning-rate selection and complete losses}
\label{app:appendix_lr}

The short-schedule sweep uses 320 updates and a full decay schedule compressed into that horizon. Its preferred rates are much larger than those supported by 5,120-step training (\Cref{tab:lr_cliff,tab:controls}). This makes it a test of a different schedule, rather than a cheap estimate of the full-budget optimum.

\begin{table}[tbp]\centering\footnotesize
\setlength{\tabcolsep}{5.2pt}
\caption{\textbf{Short-schedule (320-step) versus full-budget (5,120-step) matrix learning-rate sweeps at 130M (seed 0).} Entries are final held-out losses $\Lsel$ in nats; bold marks each row's minimum within that horizon. All four core full-budget curves achieve an interior minimum at $0.0160$, whereas the compressed 320-step schedule shifts apparent optima $4\text{--}6\times$ higher.}\label{tab:lr_cliff}
\begin{tabularx}{\textwidth}{@{}>{\raggedright\arraybackslash}Xrrrrrrrr@{}}
\toprule
Method & 0.0113 & 0.0160 & 0.0226 & 0.0320 & 0.0453 & 0.0640 & 0.0905 & 0.1280 \\
\midrule
\addlinespace[2pt]\multicolumn{9}{@{}l}{\textit{Short-schedule proxy sweep: 320 steps (0.17B tokens, seed 0)}}\\
MuonT & --- & 4.5661 & 4.4221 & 4.3347 & 4.2720 & \textbf{4.2510} & 4.2533 & 4.3021 \\
OD-All & --- & 4.5787 & 4.4312 & 4.3326 & 4.2577 & 4.2339 & \textbf{4.2314} & 4.2693 \\
OD-Outer & --- & 4.5854 & 4.4417 & 4.3337 & 4.2702 & 4.2442 & \textbf{4.2367} & 4.2798 \\
OD-Outer-F & --- & 4.5887 & 4.4412 & 4.3346 & 4.2748 & \textbf{4.2411} & 4.2436 & 4.2878 \\
\addlinespace[5pt]\multicolumn{9}{@{}l}{\textit{Full-budget verification: 5,120 steps (2.68B tokens, seed 0)}}\\
Muon & 3.2496 & \textbf{3.2461} & 3.2524 & 3.2667 & --- & --- & --- & --- \\
MuonT & 3.2404 & \textbf{3.2342} & 3.2354 & 3.2398 & --- & 3.2775 & --- & --- \\
OD-All & 3.2383 & \textbf{3.2318} & 3.2333 & 3.2403 & --- & --- & 3.3200 & --- \\
OD-Outer-F & --- & --- & \textbf{3.2318} & --- & --- & 3.2752 & --- & --- \\
\rowcolor{rowtint}[0pt][0pt]\textbf{WT-MLP-Lag} & 3.2363 & \textbf{3.2303} & 3.2311 & 3.2356 & --- & --- & --- & --- \\
\bottomrule\end{tabularx}
\end{table}


Comparing the two blocks of \Cref{tab:lr_cliff} directly illustrates the transfer gap: rates of $0.064\text{--}0.0905$ win the 320-step sweep, whereas full-budget losses achieve a strict interior minimum at $0.0160$ across all four core methods (rising at both $0.0113$ and $0.0226\text{--}0.0320$) and degrade by $0.04\text{--}0.09$ nats at the short-schedule choices. These checked transfers establish the short-horizon mismatch without assuming every unsampled rate is unstable.

At 130M and 80 TPP, the unmatched hotter WT-MLP-Lag run finishes with a Paloma macro difference of $+0.00358$ nats from MuonT; the matched-rate run gives $-0.00566$ (\Cref{tab:absolute_results}). Learning-rate transfer can therefore change the sign of the apparent output benefit. Both the matched-optimum ($\eta=0.01600$) and hotter ($\eta=0.02263$) fresh-seed replicates at 20 TPP are recorded in \Cref{tab:absolute_results,tab:seeds}.

\begin{table}[tbp]\centering\scriptsize
\setlength{\tabcolsep}{8.5pt}
\caption{\textbf{Complete absolute losses and learning rates across all confirmatory regimes.} All losses are cross-entropy in nats; Late-11 Held-out averages the 11 second-half evaluation checkpoints. Hotter-LR rows ($\eta=0.02263$ at 20 TPP; $\eta=0.01493$ at 80 TPP) document the learning-rate sensitivity discussed in \Cref{sec:phase2,sec:phase3}.}\label{tab:absolute_results}
\begin{tabularx}{\textwidth}{@{}>{\raggedright\arraybackslash}Xrrrrr@{}}
\toprule
Method & Matrix $\eta$ & Final Held-out & Late-11 Held-out & Final C4 & Paloma Macro \\
\midrule
\addlinespace[2pt]\multicolumn{6}{@{}l}{\textit{130M seed-0 controls: 20 TPP, 2.68B tokens (complete learning-rate bracket)}}\\
Muon & 0.01131 & 3.2496 & 3.3692 & 3.4661 & 3.4179 \\
MuonT & 0.01131 & 3.2404 & 3.3521 & 3.4589 & 3.4154 \\
OD-All & 0.01131 & 3.2383 & 3.3500 & 3.4569 & 3.4115 \\
WT-MLP-Lag & 0.01131 & 3.2363 & 3.3476 & 3.4548 & 3.4088 \\
Muon & 0.01600 & 3.2461 & 3.3824 & 3.4618 & --- \\
MuonT & 0.01600 & 3.2342 & 3.3591 & 3.4511 & --- \\
OD-All & 0.01600 & 3.2318 & 3.3555 & 3.4475 & 3.4048 \\
WT-MLP-Lag & 0.01600 & 3.2303 & 3.3544 & 3.4472 & --- \\
Muon & 0.02263 & 3.2524 & 3.4064 & 3.4675 & --- \\
MuonT & 0.02263 & 3.2354 & 3.3763 & 3.4511 & --- \\
Null-WT & 0.02263 & 3.2351 & 3.3755 & 3.4506 & --- \\
OD-All & 0.02263 & 3.2333 & 3.3724 & 3.4489 & --- \\
WT-MLP-Lag & 0.02263 & 3.2311 & 3.3716 & 3.4491 & --- \\
Muon & 0.03200 & 3.2667 & 3.4426 & 3.4819 & 3.4313 \\
MuonT & 0.03200 & 3.2398 & 3.3970 & 3.4563 & --- \\
OD-All & 0.03200 & 3.2403 & 3.3981 & 3.4563 & --- \\
WT-MLP-Lag & 0.03200 & 3.2356 & 3.3947 & 3.4533 & --- \\
\addlinespace[4pt]\multicolumn{6}{@{}l}{\textit{130M confirmation: 20 TPP, 2.68B tokens (mean of fresh seeds 1--3; matched $\eta^*$ \& hotter LR)}}\\
Muon & 0.01600 & 3.2035 & 3.3375 & 3.4640 & 3.4164 \\
MuonT & 0.01600 & 3.1916 & 3.3143 & 3.4523 & 3.4049 \\
OD-All & 0.01600 & 3.1919 & 3.3146 & 3.4514 & 3.4036 \\
\rowcolor{rowtint}[0pt][0pt]WT-MLP-Lag & 0.01600 & 3.1881 & 3.3107 & 3.4484 & 3.4005 \\
OD-All (hotter LR) & 0.02263 & 3.1911 & 3.3282 & 3.4493 & 3.4040 \\
WT-MLP-Lag (hotter LR) & 0.02263 & 3.1883 & 3.3265 & 3.4488 & 3.4009 \\
\addlinespace[4pt]\multicolumn{6}{@{}l}{\textit{130M longer training: 80 TPP, 10.74B tokens, seed 0 (matched \& hotter rates)}}\\
Muon & 0.01056 & 3.0778 & 3.2236 & 3.2989 & 3.2479 \\
MuonT & 0.01056 & 3.0675 & 3.2013 & 3.2886 & 3.2368 \\
OD-All & 0.01056 & 3.0656 & 3.1997 & 3.2861 & 3.2342 \\
\rowcolor{rowtint}[0pt][0pt]WT-MLP-Lag & 0.01056 & 3.0636 & 3.1967 & 3.2848 & 3.2312 \\
WT-MLP-Lag (hotter LR) & 0.01493 & 3.0696 & 3.2203 & 3.2900 & 3.2404 \\
\addlinespace[4pt]\multicolumn{6}{@{}l}{\textit{520M width transfer: 20 TPP, 10.74B tokens, seed 0}}\\
Muon & 0.01600 & 2.8528 & 2.9858 & 3.0769 & 3.0291 \\
MuonT & 0.02263 & 2.8463 & 2.9839 & 3.0704 & 3.0256 \\
OD-All & 0.02263 & 2.8469 & 2.9845 & 3.0708 & 3.0261 \\
\rowcolor{rowtint}[0pt][0pt]WT-MLP-Lag & 0.02263 & 2.8435 & 2.9809 & 3.0682 & 3.0202 \\
\bottomrule\end{tabularx}
\end{table}

\begin{table}[tbp]\centering\footnotesize
\setlength{\tabcolsep}{7.5pt}
\caption{\textbf{Individual fresh-seed replication at 130M (20 TPP, 2.68B tokens).} Subheaders report each seed's absolute MuonT losses ($\eta=0.01600$); rows report paired differences from same-seed MuonT in millinats (negative favors the row method; Late-11 is omitted for $\eta=0.02263$ rows because intermediate decay schedules differ). WT-MLP-Lag improves over both Muon and MuonT on every seed across all evaluation suites at both learning rates.}\label{tab:seeds}
\begin{tabularx}{\textwidth}{@{}>{\raggedright\arraybackslash}Xrrrrr@{}}
\toprule
Method & Matrix $\eta$ & $\Delta$ Held-out & Late-11 $\Delta$ & $\Delta$ C4 & $\Delta$ Paloma \\
\midrule
\addlinespace[2pt]\multicolumn{6}{@{}l}{\textit{Seed 1: MuonT reference ($\Lsel=3.1879$, Late-11 $=3.3113$, $\Lc=3.4528$, $\Lpal=3.4046$ nats)}}\\
Muon & 0.01600 & $+12.44$ & $+24.87$ & $+12.66$ & $+11.64$ \\
OD-All & 0.01600 & $+0.61$ & $+0.79$ & $-0.45$ & $-0.26$ \\
\rowcolor{rowtint}[0pt][0pt]\textbf{WT-MLP-Lag} & 0.01600 & $\mathbf{-1.58}$ & $\mathbf{-1.63}$ & $\mathbf{-3.53}$ & $\mathbf{-4.05}$ \\
OD-All (hotter LR) & 0.02263 & $-1.02$ & --- & $-4.08$ & $-2.62$ \\
WT-MLP-Lag (hotter LR) & 0.02263 & $-2.87$ & --- & $-4.35$ & $-4.35$ \\
\addlinespace[4pt]\multicolumn{6}{@{}l}{\textit{Seed 2: MuonT reference ($\Lsel=3.1634$, Late-11 $=3.2849$, $\Lc=3.4519$, $\Lpal=3.4026$ nats)}}\\
Muon & 0.01600 & $+12.97$ & $+23.99$ & $+11.98$ & $+15.59$ \\
OD-All & 0.01600 & $+0.39$ & $+0.37$ & $-1.57$ & $+0.60$ \\
\rowcolor{rowtint}[0pt][0pt]\textbf{WT-MLP-Lag} & 0.01600 & $\mathbf{-4.95}$ & $\mathbf{-5.23}$ & $\mathbf{-5.13}$ & $\mathbf{-2.09}$ \\
OD-All (hotter LR) & 0.02263 & $-0.02$ & --- & $-2.60$ & $+4.73$ \\
WT-MLP-Lag (hotter LR) & 0.02263 & $-2.23$ & --- & $-3.45$ & $-1.15$ \\
\addlinespace[4pt]\multicolumn{6}{@{}l}{\textit{Seed 3: MuonT reference ($\Lsel=3.2236$, Late-11 $=3.3466$, $\Lc=3.4520$, $\Lpal=3.4076$ nats)}}\\
Muon & 0.01600 & $+10.33$ & $+20.92$ & $+10.73$ & $+7.25$ \\
OD-All & 0.01600 & $-0.12$ & $-0.12$ & $-0.52$ & $-4.24$ \\
\rowcolor{rowtint}[0pt][0pt]\textbf{WT-MLP-Lag} & 0.01600 & $\mathbf{-3.95}$ & $\mathbf{-3.71}$ & $\mathbf{-2.90}$ & $\mathbf{-7.10}$ \\
OD-All (hotter LR) & 0.02263 & $-0.59$ & --- & $-2.04$ & $-4.77$ \\
WT-MLP-Lag (hotter LR) & 0.02263 & $-5.00$ & --- & $-2.69$ & $-6.53$ \\
\bottomrule\end{tabularx}
\end{table}


\FloatBarrier
\section{Placement ablations}
\label{app:ablations}

The ablation sweep changes where moments are collected and which factors are applied (\Cref{tab:ablations_full}). Each row retains its cohort's recorded MuonT reference. Earlier waves and the later fast-lag implementation have different numerical paths, so comparisons between those waves are exploratory. The fresh-seed confirmation and matched-path controls provide the stronger test of the final method.

\begin{table}[tbp]\centering\footnotesize
\setlength{\tabcolsep}{3.6pt}
\caption{\textbf{Complete 19-arm placement and operator ledger at 130M (seed 0, 2.68B tokens).} Cohort identifies the implementation wave providing each arm's paired MuonT reference (W0--W4, W10, or FastLag); loss columns are millinats from that reference (Late averages 11 second-half checkpoints; Train averages the final 537M tokens). Step $\Delta t$ is median overhead from the unified 8-GPU benchmark. Because earlier waves precede the unified confirmation runner, cross-wave differences are exploratory.}\label{tab:ablations_full}
\begin{tabularx}{\textwidth}{@{}>{\raggedright\arraybackslash}Xlllrrrrr@{}}
\toprule
& & ($\dm/\dff$) & ($\dm/\dff$) & \multicolumn{4}{c}{$\Delta$ vs.\ cohort MuonT (mnats)} & \\
\cmidrule(lr){5-8}
Method & Cohort & Inputs & Outputs & Final & Late & Late C4 & Train & Step $\Delta t$ \\
\midrule
\addlinespace[2pt]\multicolumn{9}{@{}l}{\textit{Reference baselines and execution-path control}}\\
Muon & W0 & $I$ / $I$ & $I$ / $I$ & $+9.48$ & $+5.50$ & $+5.34$ & $+6.16$ & $-1.4\%$ \\
MuonT & W0 & Dense / Dense & $I$ / $I$ & $+0.00$ & $+0.00$ & $+0.00$ & $+0.00$ & $+0.0\%$ \\
Null-WT & W3 & Dense / Dense & $I$ / $I$ (solved) & $-1.33$ & $-0.65$ & $-2.04$ & $-0.69$ & $+12.3\%$ \\
Null-Dense & W1 & Dense / Dense & $I$ / $I$ (solved) & $-1.17$ & $-1.29$ & $-2.30$ & $-1.36$ & $+25.7\%$ \\
Null-OD & W10 & Dense / Dense & $I$ / $I$ (solved) & $-1.40$ & $-2.01$ & $-1.22$ & $-1.14$ & $+1.4\%$ \\
\addlinespace[4pt]\multicolumn{9}{@{}l}{\textit{Replacing or weakening input whitening}}\\
Abl-BothDiag & W1 & Diag / Diag & Diag / Diag & $+7.33$ & $+18.01$ & $+17.69$ & $+13.37$ & $-0.4\%$ \\
Abl-OutOnly & W1 & $I$ / $I$ & Dense / Dense & $+3.25$ & $+8.71$ & $+10.90$ & $+7.87$ & $+29.8\%$ \\
Abl-DownSwap & W0 & Dense / $I$ & $I$ (\texttt{down} Dense) & $+1.06$ & $+2.51$ & $+3.61$ & $+3.18$ & $+0.3\%$ \\
WT-Sym & FastLag & Dense / Diag & Dense / Diag & $-2.99$ & $-2.52$ & $-2.03$ & $-1.26$ & $+9.3\%$ \\
WT-Sym-Lag & FastLag & Dense / Diag & Dense / Diag & $-2.46$ & $-1.13$ & $-1.60$ & $-1.48$ & $+3.3\%$ \\
\addlinespace[4pt]\multicolumn{9}{@{}l}{\textit{Output projection role and inner/outer exponent placement}}\\
Abl-DownOnly & W0 & Dense / Dense & $\hat L_{\rm down}$ only & $-1.27$ & $-1.37$ & $-0.78$ & $-1.16$ & $+2.1\%$ \\
OD-Outer & W10 & Dense / Dense & $\hat L$ Diag / $\hat L$ Diag & $-1.79$ & $-2.45$ & $-2.60$ & $-1.38$ & $+1.5\%$ \\
Abl-No2048 & W2 & Dense / Dense & Dense / $I$ & $-3.11$ & $-3.07$ & $-2.63$ & $-2.13$ & $+12.0\%$ \\
Abl-OuterDense & W2 & Dense / Dense & $\hat L$ Dense ($L=I$) & $-4.18$ & $-3.98$ & $-4.77$ & $-3.88$ & $+26.4\%$ \\
OD-All & W2 & Dense / Dense & Diag / Diag & $-4.13$ & $-4.50$ & $-3.61$ & $-3.69$ & $+2.2\%$ \\
\addlinespace[4pt]\multicolumn{9}{@{}l}{\textit{Dense two-sided ceiling versus width thresholding and one-step lag}}\\
DT-All & W1 & Dense / Dense & Dense / Dense & $-6.14$ & $-5.78$ & $-6.38$ & $-4.97$ & $+21.6\%$ \\
DT-All-Lag & W3 & Dense / Dense & Dense / Dense & $-6.21$ & $-5.89$ & $-5.86$ & $-5.73$ & $+16.5\%$ \\
WT-MLP & W2 & Dense / Dense & Dense / Diag & $\mathbf{-6.89}$ & $\mathbf{-6.58}$ & $-6.08$ & $-5.06$ & $+11.3\%$ \\
\rowcolor{rowtint}[0pt][0pt]\textbf{WT-MLP-Lag} & W4 & Dense / Dense & Dense / Diag & $\mathbf{-6.56}$ & $\mathbf{-5.68}$ & $\mathbf{-4.99}$ & $\mathbf{-4.95}$ & $+5.3\%$ \\
\bottomrule\end{tabularx}
\end{table}


\FloatBarrier
\section{Systems measurements and domain evaluation}
\label{app:appendix_systems}

The 130M timing benchmark initializes every arm from the same step-2,500 model weights with fresh optimizer and moment state, warms up for 36 steps, and records steps 37--96. It measures rank-0 real-data training steps including data transfer and distributed gradient synchronization. It excludes compilation, evaluation, and checkpoint I/O. One 60-step window per arm cannot resolve precise percent-level throughput rankings, and the records do not provide per-rank maximum durations.

\begin{table}[tbp]\centering\footnotesize
\setlength{\tabcolsep}{9.5pt}
\caption{\textbf{Complete 19-arm 130M timing and memory benchmark on 8$\times$A100.} Sorted by median step duration over steps 37--96 (after 36 warmup steps) from identical step-2,500 weights and fresh optimizer/moment state. Both median and mean overheads ($\Delta$) are relative to MuonT.}\label{tab:bench19}
\begin{tabularx}{\textwidth}{@{}>{\raggedright\arraybackslash}Xrrrrr@{}}
\toprule
Method & Median (s) & Median $\Delta$ & Mean (s) & Mean $\Delta$ & Peak GiB \\
\midrule
Muon & 1.0810 & $-1.4\%$ & 1.0792 & $-1.6\%$ & 18.97 \\
Abl-BothDiag & 1.0926 & $-0.4\%$ & 1.0911 & $-0.5\%$ & 19.55 \\
MuonT & 1.0967 & $+0.0\%$ & 1.0971 & $+0.0\%$ & 20.54 \\
Abl-DownSwap & 1.1005 & $+0.3\%$ & 1.0976 & $+0.0\%$ & 19.41 \\
Null-OD & 1.1121 & $+1.4\%$ & 1.1121 & $+1.4\%$ & 20.55 \\
OD-Outer & 1.1131 & $+1.5\%$ & 1.1134 & $+1.5\%$ & 20.55 \\
Abl-DownOnly & 1.1202 & $+2.1\%$ & 1.1283 & $+2.8\%$ & 20.76 \\
OD-All & 1.1211 & $+2.2\%$ & 1.1196 & $+2.0\%$ & 20.55 \\
WT-Sym-Lag & 1.1331 & $+3.3\%$ & 1.1751 & $+7.1\%$ & 20.82 \\
\rowcolor{rowtint}[0pt][0pt]\textbf{WT-MLP-Lag} & 1.1548 & $+5.3\%$ & 1.1944 & $+8.9\%$ & 21.60 \\
WT-Sym & 1.1986 & $+9.3\%$ & 1.2356 & $+12.6\%$ & 20.74 \\
WT-MLP & 1.2208 & $+11.3\%$ & 1.2501 & $+13.9\%$ & 21.52 \\
Abl-No2048 & 1.2288 & $+12.0\%$ & 1.2517 & $+14.1\%$ & 21.52 \\
Null-WT & 1.2313 & $+12.3\%$ & 1.2493 & $+13.9\%$ & 21.52 \\
DT-All-Lag & 1.2782 & $+16.5\%$ & 1.4585 & $+32.9\%$ & 24.78 \\
DT-All & 1.3333 & $+21.6\%$ & 1.4595 & $+33.0\%$ & 24.21 \\
Null-Dense & 1.3785 & $+25.7\%$ & 1.4638 & $+33.4\%$ & 24.21 \\
Abl-OuterDense & 1.3861 & $+26.4\%$ & 1.5004 & $+36.8\%$ & 24.21 \\
Abl-OutOnly & 1.4232 & $+29.8\%$ & 1.4579 & $+32.9\%$ & 22.21 \\
\bottomrule\end{tabularx}
\end{table}


The 520M records distinguish common-checkpoint initialization from each arm's own warmed optimizer state. We retain the complete windows in the table rather than selecting the lowest-overhead subwindow. At eight GPUs, the matched configuration uses chunked buffers and microbatch size two. At 32 GPUs, it distributes 32 Transformer layers across four nodes. Both state initialization and rare slow steps affect cost, so the paper reports medians and means separately.

\begin{table}[tbp]\centering\footnotesize
\setlength{\tabcolsep}{9.5pt}
\caption{\textbf{Complete 520M timing and memory records across 8 and 32 A100 GPUs (1.0M batch, chunked buffers).} All rows report the complete recorded measurement window without outlier removal; own optimizer states restore each arm's converged step-5,120 moments and eigenbases. Overheads ($\Delta$) are relative to within-group MuonT.}\label{tab:bench520m}
\begin{tabularx}{\textwidth}{@{}>{\raggedright\arraybackslash}Xrrrrr@{}}
\toprule
Method & Median (s) & Median $\Delta$ & Mean (s) & Mean $\Delta$ & Peak GiB \\
\midrule
\addlinespace[2pt]\multicolumn{6}{@{}l}{\textit{8 GPUs, matched checkpoint, steps 37--96}}\\
Muon & 4.6419 & $-5.2\%$ & 4.6456 & $-5.2\%$ & 26.46 \\
MuonT & 4.8943 & $+0.0\%$ & 4.9010 & $+0.0\%$ & 29.05 \\
OD-All & 4.9863 & $+1.9\%$ & 4.9902 & $+1.8\%$ & 29.06 \\
\rowcolor{rowtint}[0pt][0pt]\textbf{WT-MLP-Lag} & 5.3219 & $+8.7\%$ & 5.3181 & $+8.5\%$ & 30.73 \\
\addlinespace[4pt]\multicolumn{6}{@{}l}{\textit{32 GPUs, matched checkpoint, steps 37--96}}\\
Muon & 1.3222 & $-4.5\%$ & 1.3241 & $-4.8\%$ & 26.46 \\
MuonT & 1.3851 & $+0.0\%$ & 1.3911 & $+0.0\%$ & 26.93 \\
OD-All & 1.4063 & $+1.5\%$ & 1.4119 & $+1.5\%$ & 26.94 \\
\rowcolor{rowtint}[0pt][0pt]\textbf{WT-MLP-Lag} & 1.4305 & $+3.3\%$ & 1.5493 & $+11.4\%$ & 27.88 \\
\addlinespace[4pt]\multicolumn{6}{@{}l}{\textit{32 GPUs, own optimizer states, steps 17--76}}\\
Muon & 1.3124 & $-5.0\%$ & 1.3163 & $-5.0\%$ & 26.46 \\
MuonT & 1.3809 & $+0.0\%$ & 1.3853 & $+0.0\%$ & 26.93 \\
OD-All & 1.4203 & $+2.9\%$ & 1.4453 & $+4.3\%$ & 26.94 \\
\rowcolor{rowtint}[0pt][0pt]\textbf{WT-MLP-Lag} & 1.4578 & $+5.6\%$ & 1.5388 & $+11.1\%$ & 27.88 \\
\addlinespace[4pt]\multicolumn{6}{@{}l}{\textit{8 GPUs, own optimizer states, steps 17--76}}\\
Muon & 4.6383 & $-4.9\%$ & 4.6396 & $-5.3\%$ & 26.46 \\
MuonT & 4.8779 & $+0.0\%$ & 4.8994 & $+0.0\%$ & 29.05 \\
OD-All & 5.0287 & $+3.1\%$ & 5.0533 & $+3.1\%$ & 29.07 \\
\rowcolor{rowtint}[0pt][0pt]\textbf{WT-MLP-Lag} & 5.6081 & $+15.0\%$ & 5.6098 & $+14.5\%$ & 30.74 \\
\bottomrule\end{tabularx}
\end{table}


Paloma scores below use identical evaluation procedures within each regime and an unweighted mean over 16 domains. The three-seed Paloma macro differences are in \Cref{tab:seeds}; the full domain table includes that three-seed mean alongside the longer and wider single-seed comparisons. Win counts do not quantify uncertainty across training seeds.

\begin{table}[tbp]\centering\small
\setlength{\tabcolsep}{7.0pt}
\caption{\textbf{Per-domain Paloma loss differences from MuonT across all 16 domains.} Entries are millinats ($\text{Method}-\text{MuonT}$; negative favors the method, bold marks the stronger output approximation within each regime). The 130M 20-TPP columns average paired differences across fresh seeds 1--3; 130M 80-TPP and 520M 20-TPP use single-seed matched-rate runs.}\label{tab:paloma_full}
\begin{tabularx}{\textwidth}{@{}>{\raggedright\arraybackslash}Xrrrrrr@{}}
\toprule
& \multicolumn{2}{c}{130M, 20 TPP (3 seeds)} & \multicolumn{2}{c}{130M, 80 TPP} & \multicolumn{2}{c}{520M, 20 TPP} \\
\cmidrule(lr){2-3}\cmidrule(lr){4-5}\cmidrule(lr){6-7}
Paloma domain & OD-All & WT-Lag & OD-All & WT-Lag & OD-All & WT-Lag \\
\midrule
C4-EN & $-0.85$ & $\mathbf{-3.86}$ & $-2.55$ & $\mathbf{-3.80}$ & $+0.39$ & $\mathbf{-2.19}$ \\
C4 domains & $-0.98$ & $\mathbf{-3.78}$ & $-1.85$ & $\mathbf{-4.07}$ & $+0.56$ & $\mathbf{-1.66}$ \\
mC4 & $+0.29$ & $\mathbf{-3.73}$ & $-1.08$ & $\mathbf{-2.66}$ & $-1.20$ & $\mathbf{-1.94}$ \\
Dolma & $-1.12$ & $\mathbf{-4.13}$ & $-2.23$ & $\mathbf{-4.35}$ & $+0.16$ & $\mathbf{-2.66}$ \\
RefinedWeb & $-1.10$ & $\mathbf{-3.96}$ & $-3.43$ & $\mathbf{-5.00}$ & $-0.79$ & $\mathbf{-3.79}$ \\
RedPajama & $-0.73$ & $\mathbf{-2.81}$ & $-3.63$ & $\mathbf{-4.11}$ & $-1.33$ & $\mathbf{-2.75}$ \\
Wikipedia & $-0.75$ & $\mathbf{-3.88}$ & $-2.51$ & $\mathbf{-4.80}$ & $+3.18$ & $\mathbf{-1.79}$ \\
S2ORC & $-0.51$ & $\mathbf{-4.35}$ & $-3.11$ & $\mathbf{-4.78}$ & $-0.31$ & $\mathbf{-2.33}$ \\
Programming & $+0.51$ & $\mathbf{-4.29}$ & $-0.77$ & $\mathbf{-0.91}$ & $+0.56$ & $\mathbf{-1.63}$ \\
WikiText-103 & $\mathbf{+2.05}$ & $+2.84$ & $-6.14$ & $\mathbf{-11.00}$ & $\mathbf{-4.46}$ & $-4.06$ \\
Penn Treebank & $\mathbf{-9.42}$ & $-1.52$ & $-15.72$ & $\mathbf{-19.42}$ & $+1.29$ & $\mathbf{-20.91}$ \\
Subreddits & $+0.62$ & $\mathbf{-1.85}$ & $\mathbf{-3.34}$ & $-3.03$ & $+1.45$ & $\mathbf{-0.63}$ \\
TwitterAAE & $-6.94$ & $\mathbf{-19.72}$ & $+18.06$ & $\mathbf{-2.21}$ & $+1.31$ & $\mathbf{-40.73}$ \\
Manosphere & $-1.50$ & $\mathbf{-7.05}$ & $-0.57$ & $\mathbf{-1.46}$ & $\mathbf{-4.68}$ & $-2.83$ \\
4chan & $+2.33$ & $\mathbf{-3.40}$ & $-7.74$ & $\mathbf{-9.18}$ & $+0.18$ & $\mathbf{-0.01}$ \\
Gab & $-2.62$ & $\mathbf{-5.16}$ & $-4.86$ & $\mathbf{-9.87}$ & $+11.19$ & $\mathbf{+2.83}$ \\
\midrule
\rowcolor{rowtint}[0pt][0pt]\textbf{16-domain macro mean} & $-1.30$ & $\mathbf{-4.41}$ & $-2.59$ & $\mathbf{-5.66}$ & $+0.47$ & $\mathbf{-5.44}$ \\
\textbf{Domains won vs.\ MuonT} & 11/16 & \textbf{15/16} & 15/16 & \textbf{16/16} & 6/16 & \textbf{15/16} \\
\bottomrule\end{tabularx}
\end{table}

